\begin{resumo}
Modelos de aprendizado de máquina em produção degradam ao longo do tempo
devido ao data drift, mudança na distribuição dos dados que compromete
silenciosamente a validade do modelo treinado. Embora a detecção desse
fenômeno seja relativamente madura, a escolha de como reagir a um drift
detectado permanece pouco investigada de forma sistemática, sobretudo quanto
à interação entre o padrão temporal do drift, sua severidade e o
classificador monitorado. Este trabalho propõe um experimento fatorial
controlado, em ambiente de aprendizado em stream, para avaliar a eficácia de
quatro estratégias de retreinamento (reset completo, janela deslizante,
retreinamento incremental e reativação de modelos históricos), além de uma
condição de controle sem retreinamento, sob quatro padrões temporais de
drift, três níveis de severidade e três classificadores com mecânicas
distintas (Naive Bayes, Random Forest e XGBoost). A severidade é
operacionalizada por uma medida de distância de Hellinger com forma fechada,
calculada tanto a partir das distribuições geradoras conhecidas quanto
estimada pelo detector adotado (HDDDM), permitindo validar essa estimativa
contra o valor real. Como desdobramento exploratório, investiga-se ainda se
sinais observáveis do detector permitem esboçar uma orientação preliminar de
escolha de estratégia aplicável a pipelines de MLOps, nos quais o tipo real
de drift nunca é diretamente conhecido.
\vspace{\onelineskip}
    
\noindent \textbf{Palavras-chave}: data drift. concept drift. aprendizado em
\end{resumo}
